\documentclass[fleqn,10pt]{wlscirep}
\usepackage[utf8]{inputenc}
\usepackage[T1]{fontenc}
\title{Multi-Domain Deepfake Detection via Spatial-Frequency Feature Learning and Bidirectional Multi-Head Attention Fusion}

\author[1, *]{Abdelrahman Hassan }
\author[2]{Ahmed Hassan }
\author[3]{Hossam Abdelmunim }
\affil[1, 2, 3]{Computer and Systems Engineering Department, Faculty of Engineering, Ain Shams University, Cairo, Egypt}

\affil[*]{2201662@eng.asu.edu.eg }


\keywords{Deepfake Detection, Generative Adversarial Networks, Diffusion Models, Spatial-Domain Analysis, Wavelet Frequency Features, Multi-Head Attention, Feature Fusion, Multimedia Forensics, ResNet Backbone, Synthetic Image Classification, FaceForensics++, CelebA Dataset, Zenodo Dataset, t-SNE Visualization, Grad-CAM Explainability, Frequency-Domain Forensics, Image Authenticity Verification}

\begin{abstract}
The rapid advancement of generative modeling has made synthetic facial media detection a critical challenge in multimedia forensics. As deepfakes synthesized by Generative Adversarial Networks (GANs) and modern denoising diffusion models achieve high visual fidelity, single-domain detectors often struggle to simultaneously capture localized spatial blending anomalies and global spectral irregularities within heterogeneous multi-generator benchmarks. This paper presents a modular multi-domain deepfake detection framework that integrates spatial-domain and frequency-domain representations through a bidirectional Multi-Head Attention (MHA) fusion architecture. To evaluate the framework under controlled yet diverse generation paradigms, we curate a large-scale multi-source benchmark of 1,016,855 facial images---comprising 263,953 authentic and 752,902 synthetic samples---aggregated from CelebA, FaceForensics++, FakeAVCeleb, Celeb-DF, and Zenodo diffusion repositories (ADM, DDPM, DDIM, and DiffSwap), partitioned under strict identity-level and video-level separation. We investigate three detection configurations: (1) a modified spatial-domain detector utilizing an RGB ResNet-50 backbone with a 128-dimensional latent bottleneck, (2) a lightweight frequency-domain detector leveraging a three-level Haar wavelet packet decomposition (192 sub-bands) coupled with a convolutional network, and (3) a bidirectional MHA fusion module that dynamically weights frozen 128-dimensional spatial and frequency embeddings via multi-head cross-attention subspaces. Within the evaluated multi-generator benchmark, the spatial expert achieves 98.73\% accuracy (1.54\% EER), the wavelet expert achieves 95.42\% accuracy (5.92\% EER), and the proposed MHA fusion model achieves 98.82\% accuracy, a 0.9992 ROC-AUC, and a 1.27\% Equal Error Rate (EER). Multi-run statistical validation ($N=4$ independent seeds) confirms that while static concatenation achieves a marginally higher single-threshold accuracy (98.84\% at $\tau=0.5$), bidirectional MHA yields statistically significant improvements ($p < 0.01$) in threshold-independent calibration (reducing $1-\text{ROC-AUC}$ from 0.0100 to 0.0008 and EER from 1.32\% to 1.27\%) while restoring real-class sensitivity (97.20\% recall vs.\ 88.07\% in the frequency-only stream). Modularity experiments with Xception and ConvNeXt backbones, embedding-dimension ablations, white-box FGSM stress testing, and t-SNE and Grad-CAM analyses further demonstrate the interpretability and within-benchmark effectiveness of dynamic spatial-frequency routing, while highlighting open challenges in zero-shot cross-dataset generalization.
\end{abstract}
\begin{document}

\flushbottom
\maketitle
\thispagestyle{empty}

\section*{Introduction}
\subsection*{Rapid Evolution of Synthetic Facial Generation Technologies}
Generative Adversarial Networks (GANs) and modern denoising diffusion models now synthesize facial imagery that routinely bypasses human visual inspection. Despite their photorealistic appearance, these generative pipelines leave subtle structural and statistical traces introduced during spatial upsampling, facial boundary blending, and iterative score-based denoising \cite{ref1, ref2, ref3}. Detecting these manipulations across both legacy GAN architectures and contemporary diffusion frameworks requires forensic models that capture complementary signal irregularities rather than relying solely on surface-level visual cues.

\subsection*{Challenges in Multi-Generator Deepfake Forensics}
Forensic detectors designed exclusively around spatial-domain pixel statistics or isolated frequency-domain spectra face inherent trade-offs when evaluated across diverse synthesis families. Spatial convolutional networks excel at identifying localized blending boundaries, color mismatches, and warping defects typical of face-swapping pipelines, yet they can be sensitive to lossy compression and spatial smoothing \cite{ref4, ref5}. Conversely, frequency-domain detectors effectively isolate periodic upsampling grid artifacts and abnormal high-frequency energy distributions across wavelet or Fourier sub-bands, but frequently exhibit elevated false-positive rates on authentic images whose high-frequency textures are altered by benign camera post-processing \cite{ref8}. Consequently, single-domain detectors often struggle to maintain balanced sensitivity and specificity when confronted with a heterogeneous mixture of GAN and diffusion manipulations.

\subsection*{Limitations of Static Multi-Domain Fusion and Scope of Study}
Although recent studies have explored hybrid spatial-frequency architectures, many existing multi-domain approaches rely on static fusion strategies, such as fixed feature concatenation or late score averaging \cite{ref6, ref7}. Static fusion applies an invariant global weighting across modalities regardless of the input sample's degradation state or generator family. When an input image undergoes spatial perturbation or exhibits domain-specific conflicts---such as clean spatial semantics paired with noisy high-frequency spectra---static concatenation cannot dynamically suppress the unreliable stream.

To address this limitation within a unified multi-generator benchmark, this study proposes a modular, two-stage multi-domain detection framework based on frozen domain experts and a trainable bidirectional Multi-Head Attention (MHA) fusion core. Rather than optimizing a monolithic end-to-end multi-branch network that is prone to cross-modality co-adaptation and overfitting, our design decouples representation learning from cross-domain decision routing. Specifically, our contributions are structured as follows:
\begin{enumerate}
    \item \textbf{Decoupled Expert Representation Learning with Compact Bottlenecks:} We adapt and retrain a spatial RGB ResNet-50 expert (modified from Wang et al. \cite{ref4}) and a lightweight three-level Haar wavelet packet convolutional expert (modified from Wolter et al. \cite{ref8}, requiring only 1.73M parameters across 192 frequency sub-bands), equipping each with a supervised 128-dimensional latent bottleneck prior to freezing their weights.
    \item \textbf{Bidirectional Multi-Head Cross-Attention Fusion:} We formulate a mathematical cross-attention mechanism in which spatial embeddings query frequency keys/values and frequency embeddings reciprocally query spatial representations across multiple parallel attention subspaces ($N_{\text{heads}}=4$). We provide a complete layer-by-layer architectural specification, computational complexity analysis (FLOPs, memory footprint, and latency), and multi-seed statistical variance evaluation ($N=4$) comparing dynamic MHA routing against static concatenation (\texttt{Fusion\_128}).
    \item \textbf{Large-Scale Multi-Generator Benchmark and Modular Extensibility:} We evaluate the framework on a curated benchmark of 1,016,855 images spanning GAN and diffusion manipulations (FaceForensics++, CelebA, FakeAVCeleb, Celeb-DF, and Zenodo diffusion subsets) under strict video-level and identity-level splits. We further demonstrate plug-and-play modularity by substituting the spatial stream with Xception and ConvNeXt backbones, and we candidly delineate the boundaries between strong within-benchmark multi-generator discrimination and the open challenge of zero-shot external cross-dataset generalization.
\end{enumerate}

\section*{Background and Related Work in Deepfake Forensics}
\subsection*{Spatial-Domain Deepfake Detection}
Spatial-domain forensic detectors operate directly on RGB pixel intensities to uncover local synthesis defects, boundary blending seams, and geometric inconsistencies. Early facial forgery detectors adopted deep convolutional backbones originally engineered for visual recognition and face verification, such as Inception-ResNet-v1 \cite{ref14} and Xception \cite{ref28}. By combining multi-scale inception modules with residual skip connections or depthwise separable convolutions, these networks mitigate vanishing gradients and learn discriminative local-to-global spatial representations. MesoNet \cite{ref24} further demonstrated that compact convolutional architectures targeting mesoscopic image properties could identify early face-swapping and facial reenactment manipulations \cite{ref18, ref20}.

A foundational milestone in generalizable spatial detection was established by Wang et al. \cite{ref4}, who demonstrated that standard ResNet-50 classifiers trained on a single CNN generator (ProGAN) with aggressive Gaussian blur and JPEG compression augmentations could detect spatial upsampling artifacts across multiple unseen CNN architectures. More recently, Ricker et al. \cite{ref5} extended spatial-domain forensics to denoising diffusion models (including ADM \cite{ref3}, DDPM \cite{ref1}, and DDIM \cite{ref2}), showing that iterative reverse-diffusion sampling distributes subtle pixel-level synthesis errors globally across the image rather than concentrating them strictly at blending boundaries. Recent spatial and vision-transformer approaches, including UCF \cite{ref31}, forgery semantics decoupling \cite{ref29}, and hybrid CNN-ViT architectures \cite{ref30}, have further refined spatial artifact isolation. Nevertheless, purely spatial detectors remain vulnerable when aggressive post-processing obscures pixel-level textures.

\subsection*{Frequency-Domain and Wavelet-Packet Forensics}
To capture periodic traces introduced by transposed convolutions, upsampling filters, and spectral regularizers, frequency-domain detectors transform inputs into Fourier, Discrete Cosine Transform (DCT), or wavelet representations. While standard Fast Fourier Transform (FFT) spectra discard spatial localization, Wolter et al. \cite{ref8} introduced full Wavelet Packet Decomposition (WPD) for deepfake detection. By recursively partitioning both low-frequency approximation and high-frequency detail components across multiple scales, WPD preserves joint spatial-frequency localization and exposes abnormal cross-band energy correlations that synthetic generators fail to match to natural camera imagery. Despite their resilience to moderate spatial perturbations, standalone wavelet detectors lack high-level semantic context and often suffer from reduced recall on authentic faces containing sharp illumination transitions or complex background textures.

\subsection*{Multi-Domain Fusion and Cross-Modal Feature Interaction}
To combine the high semantic sensitivity of RGB networks with the spectral specificity of frequency transforms, recent forensic frameworks have increasingly turned to multi-domain fusion \cite{ref6, ref7}. However, conventional multi-branch detectors typically merge extracted feature vectors via static concatenation or fixed linear projection. Such designs cannot adaptively modulate modality contributions when one stream yields ambiguous or degraded signals.

Methodologically, resolving cross-domain discrepancies in multimedia forensics parallels broader advances in multi-source feature interaction, cross-modal coherence modeling, and iterative relation reasoning across multimedia computing. Recent studies have demonstrated that when heterogeneous representations exhibit divergent noise patterns or semantic abstraction levels, structured cross-modal interaction significantly outperforms static feature aggregation. For instance, Ye et al. \cite{ref32} proposed a Dual-path Collaborative Generation Network (ACM MM 2024) that coordinates parallel representation pathways through mutual guidance, illustrating how dual-branch architectures benefit from explicit inter-path collaboration rather than isolated processing. Similarly, in modeling cross-source consistency, Ye et al. \cite{ref33} (IEEE TIP 2025) demonstrated that explicitly enforcing structural coherence across complementary representation views suppresses spurious, view-specific noise. Furthermore, Ye et al. \cite{ref34} (ACM MM 2025) introduced multi-round mutual relation extraction to iteratively refine interdependent feature pairs, while Chen et al. \cite{ref35} (IEEE TIP 2026) developed a subjective-objective correlated network that aligns heterogeneous feature spaces by modeling their intrinsic cross-source correlations.

Inspired by these methodological principles of dual-path collaboration \cite{ref32}, cross-view coherence \cite{ref33}, mutual relation reasoning \cite{ref34}, and heterogeneous feature correlation \cite{ref35}, our proposed Multi-Head Attention (MHA) fusion module treats the spatial (RGB) and frequency (wavelet packet) embeddings as interactive multi-source representations. Rather than merely concatenating static descriptors, our bidirectional cross-attention block requires the spatial embedding to query frequency subspaces for spectral corroboration while the frequency embedding simultaneously queries spatial subspaces for structural context. This mutual querying acts as a soft cross-domain consistency constraint, dynamically calibrating per-sample modality weights and suppressing conflicting forensic noise prior to classification.

\section*{Proposed Multi-Domain Detection Framework}
\subsection*{First Stream: Modified Spatial-Domain Expert (Wang2020\_128)}
\begin{figure}[ht]
\centering
\includegraphics[width=\linewidth]{fig1_Wang2020_modified_128D}
\caption{Modified RGB Wang2020 Model (Spatial Domain) with ResNet-50 Backbone and 128-D latent projection. The source facial image is an entirely synthetic, AI-generated sample sourced from ThisPersonDoesNotExist.com, used strictly for architectural illustration (it does not depict a real human individual).}
\label{fig:fig1_Wang2020_modified_128D}
\end{figure}
Our spatial-domain stream, illustrated in Figure \ref{fig:fig1_Wang2020_modified_128D}, adapts the ResNet-50 forensic detector of Wang et al. \cite{ref4, ref25} and retrains it on our multi-generator dataset. While the original architecture maps the 2048-dimensional global average pooled ResNet-50 feature vector directly to a single binary logit, we insert a compact 128-dimensional latent bottleneck layer immediately preceding the binary classification head. During Phase 1 training, this modified network is optimized end-to-end so that binary cross-entropy supervision forces the 128-dimensional bottleneck to encode highly discriminative spatial forgery cues, including boundary blending seams, color-space discontinuities, and localized warping artifacts.

Specifically, an input RGB image resized to $224 \times 224 \times 3$ is normalized using standard ImageNet statistics and processed through the ResNet-50 stages to yield a feature vector $h_s \in \mathbb{R}^{2048}$. A linear projection maps $h_s$ to the 128-dimensional embedding space, followed by a Rectified Linear Unit (ReLU) activation and Dropout ($p = 0.5$) for regularization:
\begin{equation}
X_s = \text{Dropout}\left(\text{ReLU}\left(W_{s,\text{proj}} h_s + b_{s,\text{proj}}\right)\right) \in \mathbb{R}^{128}
\end{equation}
where $W_{s,\text{proj}} \in \mathbb{R}^{128 \times 2048}$ and $b_{s,\text{proj}} \in \mathbb{R}^{128}$. During Phase 1 training, a temporary linear head computes the binary logit using Binary Cross-Entropy with Logits Loss (\texttt{BCEWithLogitsLoss}) and the Adam optimizer. Once Phase 1 converges, the temporary classification head is discarded and all parameters of the spatial expert are frozen, yielding a deterministic spatial representation extractor $X_s$ for Phase 2 fusion.

\subsection*{Second Stream: Modified Frequency-Domain Expert (Wolter2021\_128)}
\begin{figure}[ht]
\centering
\includegraphics[width=\linewidth]{fig2_WolterWaveletesWolter2021_modified_128D.jpg}
\caption{Modified Wavelet Packet Wolter2021 Architecture (Frequency Domain) with 3-level Haar decomposition (192 sub-bands) and 128-D latent projection. The source facial image is an entirely synthetic, AI-generated sample sourced from ThisPersonDoesNotExist.com, used strictly for architectural illustration (it does not depict a real human individual).}
\label{fig:fig2_WolterWaveletesWolter2021_modified_128D}
\end{figure}
Our frequency-domain stream, shown in Figure \ref{fig:fig2_WolterWaveletesWolter2021_modified_128D}, builds upon the Wavelet Packet CNN architecture introduced by Wolter et al. \cite{ref8}. Rather than operating on raw spatial pixels, this branch applies a three-level 2D Wavelet Packet Decomposition (WPD) to each RGB channel, partitioning the image into $4^3 = 64$ frequency sub-bands per color channel ($3 \times 64 = 192$ total spectral channels) that capture approximation, horizontal, vertical, and diagonal oscillations across multiple frequency scales.

The selection of a three-level decomposition ($L=3$, 192 channels) reflects a principled trade-off between spectral resolution and representation complexity \cite{ref8}. A shallower two-level decomposition ($L=2$, yielding $3 \times 4^2 = 48$ channels) provides coarse frequency partitioning that fails to isolate narrowband periodic upsampling artifacts and subtle high-frequency diffusion denoising residuals. Conversely, a four-level decomposition ($L=4$, yielding $3 \times 4^4 = 768$ channels at reduced spatial resolution) quadruples the input channel dimensionality, substantially increasing computational overhead and promoting overfitting to camera-specific sensor noise. The three-level representation ($192 \times 28 \times 28$) preserves sufficient spatial-frequency granularity to isolate synthesis traces while keeping the downstream convolutional network compact. To stabilize heavy-tailed wavelet coefficient distributions and accentuate subtle high-frequency energy deviations, the magnitude of the wavelet packet coefficients $C_{i,j,k}$ is log-scaled prior to convolutional processing:
\begin{equation}
\tilde{C}_{i,j,k} = \log\left(\lvert C_{i,j,k} \rvert + \epsilon\right), \quad \epsilon = 10^{-6}
\end{equation}

\subsubsection*{Wavelet Expert Architectural Specifications}
To ensure full reproducibility, the exact parameterization of the \texttt{Wolter2021\_128} expert is specified as follows. The three-level WPD is computed using the orthogonal Haar wavelet family (\texttt{db1}) with reflection boundary padding to prevent artificial edge ringing across the 192 spectral sub-bands. The resulting $192$-channel tensor is processed by a lightweight Convolutional Neural Network (CNN) comprising sequential $3 \times 3$ convolutional blocks with zero-padding ($\text{padding}=1$), two-dimensional Batch Normalization (\texttt{BatchNorm2d}), ReLU activations, and $2 \times 2$ max-pooling layers, followed by global average pooling to form a 512-dimensional spectral descriptor $h_f \in \mathbb{R}^{512}$.

Analogous to the spatial expert, we introduce a 128-dimensional latent bottleneck layer ($\text{Linear}(512 \to 128) \to \text{ReLU} \to \text{Dropout}(0.5)$) prior to the binary classification head. After end-to-end Phase 1 training on the spectral stream, the binary classification head is removed and all wavelet expert weights are frozen, providing a compact 128-dimensional frequency embedding $X_f \in \mathbb{R}^{128}$ that complements the spatial stream within the evaluated benchmark.

\subsection*{Third Stream: Bidirectional Multi-Head Attention Fusion Module (MHA\_128)}
\begin{figure}[htbp]
\centering
\includegraphics[width=0.9\linewidth]{fig3_WholeArch}
\caption{Complete architecture of the proposed multi-domain deepfake detection framework. Pretrained and frozen Wang2020 (RGB) and Wolter2021 (Wavelet) experts extract 128-dimensional domain embeddings, which interact through parallel 4-head cross-attention blocks, residual LayerNorm, and feed-forward networks prior to classification. The source facial image is an entirely synthetic, AI-generated sample sourced from ThisPersonDoesNotExist.com, used strictly for architectural illustration (it does not depict a real human individual).}
\label{fig:fig3_WholeArch}
\end{figure}
Our central fusion module, depicted in Figure \ref{fig:fig3_WholeArch}, is a bidirectional Multi-Head Attention (MHA) network that integrates the frozen 128-dimensional embeddings extracted from the spatial expert ($X_s \in \mathbb{R}^{1 \times D}$) and the frequency expert ($X_f \in \mathbb{R}^{1 \times D}$), where $D = 128$.

To clarify how multi-head attention operates meaningfully on compact domain embeddings, each 128-dimensional vector is decomposed across $N_{\text{heads}} = 4$ parallel subspace heads of dimension $D_k = D / N_{\text{heads}} = 32$. Rather than computing a single scalar weight between two monolithic vectors, the four attention heads project $X_s$ and $X_f$ into distinct 32-dimensional semantic and spectral subspaces, allowing the module to evaluate fine-grained subspace correlations across domains. Specifically, for each head $h \in \{1, \dots, N_{\text{heads}}\}$ in the two parallel cross-attention blocks (spatial querying frequency $s \to f$, and frequency querying spatial $f \to s$), the head-specific Query, Key, and Value projections are computed as:
\begin{align}
Q_s^{(h)} = X_s W_s^{Q,(h)}, \quad K_f^{(h)} = X_f W_f^{K,(h)}, \quad V_f^{(h)} = X_f W_f^{V,(h)} \\
Q_f^{(h)} = X_f W_f^{Q,(h)}, \quad K_s^{(h)} = X_s W_s^{K,(h)}, \quad V_s^{(h)} = X_s W_s^{V,(h)}
\end{align}
where $W_d^{Q,(h)}, W_d^{K,(h)}, W_d^{V,(h)} \in \mathbb{R}^{D \times D_k}$ (with learnable biases) for domain $d \in \{s, f\}$. Within each head subspace, scaled dot-product cross-attention interactions are computed and concatenated across all $N_{\text{heads}} = 4$ heads via output projection matrices $W_{s \to f}^O, W_{f \to s}^O \in \mathbb{R}^{D \times D}$:
\begin{align}
A_{s \to f} &= \text{Concat}\left(\text{head}_{s \to f}^{(1)}, \dots, \text{head}_{s \to f}^{(4)}\right) W_{s \to f}^O, \quad \text{head}_{s \to f}^{(h)} = \text{Softmax}\left(\frac{Q_s^{(h)} \left(K_f^{(h)}\right)^T}{\sqrt{D_k}}\right) V_f^{(h)} \\
A_{f \to s} &= \text{Concat}\left(\text{head}_{f \to s}^{(1)}, \dots, \text{head}_{f \to s}^{(4)}\right) W_{f \to s}^O, \quad \text{head}_{f \to s}^{(h)} = \text{Softmax}\left(\frac{Q_f^{(h)} \left(K_s^{(h)}\right)^T}{\sqrt{D_k}}\right) V_s^{(h)}
\end{align}

Following the bidirectional multi-head projections, each pathway incorporates a residual skip connection and Layer Normalization ($\text{LN}$) to preserve unattenuated domain information while integrating cross-domain context:
\begin{equation}
Z^{(1)}_s = \text{LN}\left(X_s + A_{s \to f}\right), \quad Z^{(1)}_f = \text{LN}\left(X_f + A_{f \to s}\right)
\end{equation}
Each normalized representation is then transformed by a two-layer position-wise Feed-Forward Network ($\text{FFN}$) with an expanded hidden dimension $H = 512$ and dropout ($p = 0.1$), followed by a second residual connection and Layer Normalization:
\begin{align}
\text{FFN}_d\left(Z^{(1)}_d\right) &= W_{2,d} \cdot \text{ReLU}\left(W_{1,d} Z^{(1)}_d + b_{1,d}\right) + b_{2,d} \\
Z_d &= \text{LN}\left(Z^{(1)}_d + \text{FFN}_d\left(Z^{(1)}_d\right)\right), \quad d \in \{s, f\}
\end{align}
where $W_{1,d} \in \mathbb{R}^{512 \times 128}$, $W_{2,d} \in \mathbb{R}^{128 \times 512}$, $b_{1,d} \in \mathbb{R}^{512}$, and $b_{2,d} \in \mathbb{R}^{128}$. This non-linear expansion and projection across $H=512$ intermediate channels enables sample-adaptive feature gating that cannot be replicated by a static linear layer.

Finally, the refined domain tokens $Z_s$ and $Z_f$ are concatenated into $[Z_s \parallel Z_f] \in \mathbb{R}^{1 \times 256}$ and passed through a three-layer Multi-Layer Perceptron (MLP) classification head to yield the scalar forgery logit $y \in \mathbb{R}$:
\begin{align}
X_{\text{fuse}} &= \text{ReLU}\left(W_{\text{fuse}} [Z_s \parallel Z_f] + b_{\text{fuse}}\right) \\
X_{\text{cls}} &= \text{ReLU}\left(W_{\text{c1}} X_{\text{fuse}} + b_{\text{c1}}\right) \\
y &= W_{\text{c2}} X_{\text{cls}} + b_{\text{c2}}
\end{align}
where $W_{\text{fuse}} \in \mathbb{R}^{128 \times 256}$, $b_{\text{fuse}} \in \mathbb{R}^{128}$, $W_{\text{c1}} \in \mathbb{R}^{64 \times 128}$, $b_{\text{c1}} \in \mathbb{R}^{64}$, $W_{\text{c2}} \in \mathbb{R}^{1 \times 64}$, and $b_{\text{c2}} \in \mathbb{R}^1$. Table \ref{tab:mha_architecture} reports the complete layer-by-layer specification and exact trainable parameter counts for the MHA fusion module.

\begin{table}[htbp]
\centering
\begin{tabular}{|l|l|c|r|}
\hline
\textbf{Module / Layer} & \textbf{Specifications} & \textbf{Output Shape} & \textbf{Trainable Parameters} \\ \hline
Spatial Stream ($X_s$) & Wang2020-128 Backbone (Frozen) & $1 \times 128$ & — \\
Frequency Stream ($X_f$) & Wolter2021-128 Backbone (Frozen) & $1 \times 128$ & — \\ \hline
\textbf{Cross-Attention Blocks} & \multicolumn{3}{c|}{(Parallel spatial $s \to f$ and frequency $f \to s$ paths)} \\ \hline
Multihead Attention ($A$) & $N_{\text{heads}}=4$, $D=128$, $D_k=32$, $\text{bias}=\text{True}$ & $1 \times 128$ & $2 \times 66{,}048 = 132{,}096$ \\
LayerNorm 1 & $\text{features}=128$ & $1 \times 128$ & $2 \times 256 = 512$ \\
Feed-Forward Network & $\text{Linear}(128 \to 512 \to 128)$, $\text{dropout}=0.1$ & $1 \times 128$ & $2 \times 131{,}712 = 263{,}424$ \\
LayerNorm 2 & $\text{features}=128$ & $1 \times 128$ & $2 \times 256 = 512$ \\ \hline
\textbf{Fusion \& Classification} & \multicolumn{3}{c|}{(Concatenation and MLP Classification Head)} \\ \hline
Concatenation ($[Z_s \parallel Z_f]$) & Combine refined spatial and frequency tokens & $1 \times 256$ & — \\
Linear Fusion & $\text{Linear}(256 \to 128)$, $\text{ReLU}$, $\text{dropout}=0.1$ & $1 \times 128$ & $32{,}896$ \\
Classification FC1 & $\text{Linear}(128 \to 64)$, $\text{ReLU}$, $\text{dropout}=0.1$ & $1 \times 64$ & $8{,}256$ \\
Classification FC2 & $\text{Linear}(64 \to 1)$ & $1 \times 1$ & $65$ \\ \hline
\textbf{Total Trainable Core} & \multicolumn{2}{l|}{MHA Fusion Classifier Module} & \textbf{437,761} \\ \hline
\end{tabular}
\caption{\label{tab:mha_architecture}Detailed architectural specifications and layer-by-layer trainable parameter counts for the proposed Multi-Head Attention (MHA) fusion model.}
\end{table}

\subsection*{Model Complexity, Computational Footprint, and Parameter Efficiency}
Because the complete inference pipeline requires executing both the spatial expert and the frequency expert before fusing their 128-dimensional embeddings, quantifying the end-to-end computational footprint is essential for assessing practical deployability. Table \ref{tab:parameter_counts} summarizes the trainable parameter counts, FP32 model memory footprints, theoretical floating-point operations (GFLOPs per $224 \times 224$ image), and average single-image inference latency measured on an NVIDIA T4 GPU (16\,GB VRAM, batch size 1).

\begin{table}[htbp]
\centering
\begin{tabular}{|l|l|r|r|r|r|}
\hline
\textbf{Model Component} & \textbf{Architecture Base} & \textbf{Parameters} & \textbf{Size (MB)} & \textbf{GFLOPs} & \textbf{Latency (ms)} \\ \hline
Spatial Expert (\texttt{Wang2020\_128}) & ResNet-50 + 128-D Proj. & 23,770,433 & 90.68 & 4.120 & 8.4 \\
Frequency Expert (\texttt{Wolter2021\_128}) & 3-Level WPD + CNN + 128-D & 1,727,553 & 6.59 & 0.480 & 4.1 \\
Simple Fusion Head (\texttt{Fusion\_128}) & Concatenation + MLP & 107,009 & 0.41 & 0.0002 & 0.3 \\
Proposed MHA Head (\texttt{MHA\_128}) & Dual Cross-Attention + MLP & 437,761 & 1.67 & 0.0009 & 0.6 \\ \hline
\textbf{Full Dual-Stream MHA Pipeline} & \textbf{Both Experts + MHA Head} & \textbf{25,935,747} & \textbf{98.94} & \textbf{4.601} & \textbf{13.1} \\ \hline
\end{tabular}
\caption{\label{tab:parameter_counts}Computational footprint comparison across sub-model experts, fusion heads, and the complete end-to-end dual-stream MHA pipeline ($224 \times 224$ input resolution, FP32 precision on NVIDIA T4 GPU).}
\end{table}

As shown in Table \ref{tab:parameter_counts}, the modified \texttt{Wolter2021\_128} wavelet expert requires only 1.73M parameters (6.59\,MB) and 0.480 GFLOPs---representing just 7.3\% of the parameter count and 11.7\% of the FLOPs of the ResNet-50 spatial backbone. Furthermore, the bidirectional \texttt{MHA\_128} fusion head adds only 0.44M parameters (1.67\,MB) and less than 0.001 GFLOPs (0.6\,ms latency). Consequently, the complete dual-stream MHA pipeline operates at 4.601 GFLOPs and 13.1\,ms per image ($\sim 76$ FPS sequential execution, or $\sim 9.0$\,ms when the two independent expert streams are executed concurrently across parallel CUDA streams), making multi-domain fusion computationally practical without duplicating heavy spatial backbones.

\subsection*{Architectural Modularity and Multi-Token Extensibility}
While our primary benchmark pairs two domain experts (spatial RGB and wavelet packet frequency), the MHA fusion module is inherently extensible. In a static concatenation model (\texttt{Fusion\_128}), adding a third or fourth forensic stream (e.g., an Xception or ConvNeXt extractor, or a temporal optical-flow specialist) alters the input dimension of the first linear layer and requires retraining the entire classifier from scratch. Within the proposed framework, any new forensic backbone can be independently trained with a 128-dimensional projection head, frozen, and appended as an additional 128-dimensional token in the multi-head attention sequence, enabling plug-and-play multi-expert integration.

\section*{Methodology and Experimental Protocol}
\subsection*{Methodological Justification of Frozen-Expert Attention Routing}
Training a multi-branch network end-to-end from scratch often leads to modality dominance, where the faster-converging spatial ResNet-50 stream dominates gradient updates and prevents the frequency branch from learning independent spectral descriptors. To prevent cross-domain co-adaptation and representation drift, we enforce a strict two-stage decoupling strategy: domain experts are first trained independently to convergence to form specialized 128-dimensional embeddings, and their parameters are subsequently frozen during Phase 2 MHA optimization. Because the attention module cannot alter the underlying ResNet-50 or WaveletPacketCNN weights, it is forced to learn instance-dependent routing rules that reconcile discrepancies between spatial and spectral cues.

\subsection*{Staged Training Protocol and Leakage Prevention}
To eliminate any risk of validation overlap or data contamination across training phases, all models adhere to a unified, strictly partitioned protocol across the 80/20 video- and identity-level dataset split (Table \ref{tab:training_summary}):
\begin{description}
\item[Phase 1: Independent Sub-Model Representation Learning] The spatial expert (\texttt{Wang2020\_128}) and frequency expert (\texttt{Wolter2021\_128}) are trained independently end-to-end (all backbone and 128-D projection parameters unfrozen) exclusively on the 80\% Training partition (805,421 images). Model selection and early stopping are monitored on the held-out 20\% Validation partition (211,434 images).
\item[Phase 2: Frozen-Backbone MHA Fusion Training] After Phase 1 convergence, the temporary binary classification heads of both experts are removed and all expert weights (including their 128-D projection layers) are permanently frozen. The \texttt{MHA\_128} fusion core (and the baseline \texttt{Fusion\_128} concatenation head) are trained exclusively on the exact same 80\% Training partition (805,421 images), ensuring that no validation samples are ever used to update either backbone or fusion parameters.
\item[Phase 3: Held-Out Evaluation and Modularity Verification] Final quantitative metrics, multi-seed variance analyses, FGSM adversarial stress tests, and modularity experiments incorporating retrained \texttt{Xception\_128} and \texttt{ConvNeXt\_128} backbones are evaluated purely in inference mode on the held-out 20\% Validation partition (211,434 images).
\end{description}

\begin{table}[htbp]
\centering
\begin{tabular}{|l|l|l|l|}
\hline
\textbf{Phase} & \textbf{Component} & \textbf{Parameter State} & \textbf{Dataset Partition Used} \\ \hline
Phase 1 & Spatial RGB Expert (\texttt{Wang2020\_128}) & Unfrozen (End-to-End) & 80\% Train Split ($n=805{,}421$) \\ \hline
Phase 1 & Frequency Expert (\texttt{Wolter2021\_128}) & Unfrozen (End-to-End) & 80\% Train Split ($n=805{,}421$) \\ \hline
Phase 2 & Spatial \& Frequency Experts & \textbf{Strictly Frozen} & 80\% Train Split (Feature Extraction Only) \\ \hline
Phase 2 & MHA Fusion Module (\texttt{MHA\_128}) & Unfrozen (Trained) & 80\% Train Split ($n=805{,}421$) \\ \hline
Phase 3 & Complete Dual-Stream Pipeline & \textbf{Strictly Frozen} (Inference) & Held-Out 20\% Val Split ($n=211{,}434$) \\ \hline
\end{tabular}
\caption{Summary of the staged training and evaluation protocol, detailing parameter freezing states and strict partition isolation across Phase 1, Phase 2, and Phase 3.}
\label{tab:training_summary}
\end{table}

\section*{Dataset Construction and Preprocessing}
\subsection*{Multi-Source Dataset Composition and Provenance Split}
To evaluate unified detection across both GAN-based facial manipulations and diffusion-generated imagery, we compiled a multi-source dataset totaling 1,016,855 facial images: 263,953 authentic (Real) images and 752,902 synthetic (Fake) images (Table \ref{tab:dataset_splits}). Authentic facial images were primarily sourced from the CelebA dataset \cite{ref36}, supplemented by pristine real video frames from FaceForensics++ \cite{ref12}, FakeAVCeleb \cite{ref11}, and Celeb-DF \cite{ref37}. Synthetic samples encompass two major generative paradigms:
\begin{enumerate}
    \item \textbf{Diffusion-Generated Imagery (120,000 samples):} Sourced from public Zenodo diffusion forensics repositories \cite{ref9, ref10, ref21}, comprising four distinct diffusion architectures---Ablated Diffusion Model (ADM) \cite{ref3}, Denoising Diffusion Probabilistic Models (DDPM) \cite{ref1}, Denoising Diffusion Implicit Models (DDIM) \cite{ref2}, and 3D-aware masked diffusion face swapping (DiffSwap) \cite{ref16}---contributing 30,000 images per generator.
    \item \textbf{GAN and Autoencoder Facial Manipulations (632,902 samples):} Extracted from FaceForensics++ (Deepfakes, FaceSwap, Face2Face, and NeuralTextures) \cite{ref12, ref18}, FakeAVCeleb \cite{ref11}, and Celeb-DF \cite{ref37}, representing diverse identity-swapping and expression-reenactment pipelines \cite{ref13, ref17, ref19, ref20, ref22, ref23}.
\end{enumerate}

To prevent identity leakage, background memorization, or near-duplicate frame contamination between splits, dataset partitioning was executed strictly at the \textbf{source-video and subject-identity level} prior to frame extraction. Specifically, all video sequences in FaceForensics++, FakeAVCeleb, and Celeb-DF were partitioned by source video ID and actor identity (80\% training videos, 20\% validation videos), and CelebA/Zenodo samples were partitioned by unique subject identity ID. Consequently, 805,421 images (211,188 real and 594,233 fake) form the Training split, and 211,434 images (52,765 real and 158,669 fake) form the held-out Validation split, with zero identity or video overlap between partitions.

\begin{table}[htbp]
\centering
\begin{tabular}{|l|c|c|r|}
\hline
\textbf{Partition Split} & \textbf{Real Images (Authentic)} & \textbf{Fake Images (GAN + Diffusion)} & \textbf{Total Images} \\ \hline
Training ($80\%$)   & 211,188 & 594,233 & 805,421 \\
Validation ($20\%$) & 52,765  & 158,669 & 211,434 \\ \hline
\textbf{Total}      & \textbf{263,953} & \textbf{752,902} & \textbf{1,016,855} \\ \hline
\end{tabular}
\caption{\label{tab:dataset_splits}Authoritative dataset distribution across Real and Fake (synthetic) classes under strict video-level and identity-level 80/20 partitioning.}
\end{table}

\subsubsection*{Batch-Level Class Balancing Strategy}
Because the aggregated training partition exhibits an approximate 1:2.81 real-to-synthetic class ratio (211,188 real vs.\ 594,233 fake samples), naive uniform sampling would bias the classifier toward the majority synthetic class. To eliminate class-prior bias during optimization, we employ batch-level inverse-frequency weighted random sampling with replacement via PyTorch's \texttt{WeightedRandomSampler}. For each class $c \in \{\text{Real}, \text{Fake}\}$ containing $N_c$ training samples, the per-sample selection weight $w_c$ is defined as:
\begin{equation}
w_c = \frac{1}{N_c}
\end{equation}
Drawing mini-batches of size 32 with probability proportional to $w_c$ oversamples authentic faces and undersamples synthetic faces within each epoch so that every training mini-batch maintains an expected 1:1 class balance. Meanwhile, the held-out validation set ($n=211{,}434$) is evaluated without resampling to reflect its natural test distribution, and both threshold-dependent metrics (Accuracy, class-wise Precision/Recall/F1) and prevalence-invariant metrics (ROC-AUC, EER) are reported.

\subsubsection*{Frame Extraction Stride and Temporal Redundancy Control}
For video-based repositories (FaceForensics++, FakeAVCeleb, and Celeb-DF), frames were extracted at a uniform rate of 10 frames per second (FPS), followed by automated face detection, bounding-box cropping, and bilinear resizing to $224 \times 224$ pixels. While a 10 FPS sampling rate introduces correlation among adjacent frames within a single video clip, this temporal density captures transient synthesis failures that occur during rapid facial motion, such as eye blinks, phoneme transitions, and head rotations. Because all splits were enforced strictly at the video and identity level prior to frame extraction, frames from any training video never appear in the validation split, ensuring that temporal density within training videos cannot leak into validation metrics.

\subsection*{Forensic Data Augmentation Pipeline}
To discourage the experts from memorizing fragile, uncompressed pixel patterns, dynamic training-time augmentations were applied exclusively to the Training split:
\begin{description}
\item[Gaussian Blur] Applied with probability $p = 0.5$, where the kernel standard deviation $\sigma$ is sampled uniformly from $\sigma \sim \mathcal{U}(0.0, 3.0)$ to simulate spatial smoothing and defocus degradation.
\item[Dynamic JPEG Compression] Applied with probability $p = 0.5$ during training, randomly selecting between OpenCV and Pillow compression engines at discrete quality factors $Q \in \{30, 75\}$ to simulate lossy social-media transcoding. During deterministic validation preprocessing, input crops are standardized at quality factor $Q = 75$ after resizing from $256 \times 256$ to $224 \times 224$ pixels via bilinear interpolation.
\end{description}

\subsection*{Ethical, Privacy, and Demographic Fairness Considerations}
All datasets analyzed in this work (CelebA \cite{ref36}, FaceForensics++ \cite{ref12}, FakeAVCeleb \cite{ref11}, Celeb-DF \cite{ref37}, and Zenodo DiffFace \cite{ref9, ref10}) were accessed under their official academic licenses. No private user data or non-public personally identifiable information (PII) was collected. Furthermore, to respect subject privacy and copyright compliance in publication, all sample facial illustrations displayed in the architectural and explainability figures (Figures \ref{fig:fig1_Wang2020_modified_128D}, \ref{fig:fig2_WolterWaveletesWolter2021_modified_128D}, \ref{fig:fig3_WholeArch}, and \ref{fig:fig8_GradCam_Fake_Image_DiffusionModel_Detected_ADM}) depict purely synthetic, non-existent AI-generated identities produced by StyleGAN (\texttt{ThisPersonDoesNotExist.com}) or ADM diffusion synthesis, rather than recognizable human individuals.

Regarding algorithmic fairness, deepfake detectors risk exhibiting disparate false-positive or false-negative rates across demographic subgroups (age, gender, ethnicity, and skin tone) if training distributions are unbalanced. Because several constituent benchmarks (such as FaceForensics++ and Zenodo diffusion subsets) do not provide verified ground-truth demographic annotations for all 1.01M frames---and assigning automated proxy labels introduces secondary classifier bias---subgroup fairness metrics could not be reliably stratified in the current study. We explicitly highlight demographic fairness auditing on annotated benchmarks (such as FakeAVCeleb and AI Foundation fairness sets) as an essential prerequisite prior to real-world deployment.

\section*{Experimental Results and Comparative Analysis}
\subsection*{Training Convergence Analysis}
\subsubsection*{Modified Spatial Expert (Wang2020\_128) Convergence}
\begin{figure}[htbp]
\centering
\includegraphics[width=\linewidth]{fig4_Wang_Modified.jpg}
\caption{Training loss and online validation convergence profiles for the modified spatial-domain model (\texttt{Wang2020\_128}).}
\label{fig:fig4_Wang_Modified}
\end{figure}
As shown in Figure \ref{fig:fig4_Wang_Modified}, the modified \texttt{Wang2020\_128} spatial expert was trained using the Adam optimizer ($\beta_1 = 0.9$, learning rate $10^{-3}$, batch size 32, logged every 400 iterations). During online training batches, validation precision approaches 99\% and online validation accuracy stabilizes near 98.2\% (reaching 98.73\% upon full deterministic checkpoint evaluation across the complete 211,434-image validation split). The training loss converges near $10^{-3}$ with mild stochastic oscillations attributable to heavy dropout ($p=0.5$) and dynamic Gaussian blur/JPEG augmentation.

\subsubsection*{Modified Frequency Expert (Wolter2021\_128) Convergence}
\begin{figure}[htbp]
\centering
\includegraphics[width=\linewidth]{fig5_Wolter_Modified}
\caption{Training loss and online validation convergence profiles for the modified wavelet-packet frequency model (\texttt{Wolter2021\_128}).}
\label{fig:fig5_Wolter_Modified}
\end{figure}
Figure \ref{fig:fig5_Wolter_Modified} illustrates the training dynamics of the \texttt{Wolter2021\_128} wavelet packet expert, optimized via Adam ($\beta_1 = 0.9$, learning rate $10^{-4}$, batch size 32, logged every 400 iterations). Online validation precision and accuracy reach approximately 96.8\% and 94.6\% during training, converging to a deterministic validation accuracy of 95.42\% at the best saved checkpoint.

\subsubsection*{Bidirectional Multi-Head Attention Fusion (MHA\_128) Convergence}
\begin{figure}[htbp]
\centering
\includegraphics[width=\linewidth]{fig6_MHA_Modified.png}
\caption{Training loss and validation convergence profiles for the proposed Bidirectional Multi-Head Attention fusion model (\texttt{MHA\_128}). Note that in the Average Precision (AP) subplot, the vertical axis is zoomed into the narrow interval $[0.99835, 0.99845]$, so the apparent visual slope represents a negligible variation of less than $0.005\%$.}
\label{fig:fig6_MHA_Modified}
\end{figure}
As depicted in Figure \ref{fig:fig6_MHA_Modified}, the \texttt{MHA\_128} fusion core was trained on the frozen expert embeddings using Adam ($\beta_1 = 0.9$, learning rate $10^{-3}$, normal weight initialization with gain $0.02$, batch size 32). Validation accuracy rapidly converges to 98.82\% and precision reaches $\sim 99\%$, while training loss stabilizes near $10^{-4}$. Notably, in the Average Precision (AP) subplot of Figure \ref{fig:fig6_MHA_Modified}, the vertical axis is automatically scaled to an extremely narrow high-performance band ($\sim 0.9984$), meaning the visual variation across iterations represents less than $0.005\%$ fluctuation around the converged AP of 0.9984.

\subsection*{Held-Out Benchmark Evaluation of Individual and Fused Models}
\subsubsection*{Modified Spatial Expert (Wang2020\_128) Evaluation}
Table \ref{tab:wangModelModQuant} and Table \ref{tab:wangModClassWise} summarize the held-out validation performance of the modified spatial expert. Across the 211,434 validation images, the model achieves 98.20\% online milestone accuracy (and 98.73\% final deterministic evaluation accuracy in Table \ref{tab:model_comparison}), a ROC-AUC of 0.9987, an Average Precision (AP) of 0.9977, an Equal Error Rate (EER) of 1.54\%, and a macro F1-score of 0.9796. On the 52,765 real images, 51,472 are correctly classified (1,293 false positives; Real recall = 0.9755, precision = 0.9838, F1 = 0.9796). On the 158,669 synthetic images, 157,511 are correctly identified (1,158 false negatives; Fake recall = 0.9927, precision = 0.9889, F1 = 0.9908).

\begin{table}[htbp]
\centering
\begin{tabular}{|l|c|}
\hline
\textbf{Metric} & \textbf{Value} \\
\hline
Accuracy (Milestone / Final Converged) & 98.20\% / 98.73\% \\
ROC AUC & 0.9987 \\
Average Precision & 0.9977 \\
Equal Error Rate (EER) & 1.54\% \\
F1 Score (Macro) & 0.9796 \\
\hline
\end{tabular}
\caption{\label{tab:wangModelModQuant}Quantitative evaluation of the modified \texttt{Wang2020\_128} spatial expert on the held-out validation partition ($n=211{,}434$).}
\end{table}

\begin{table}[htbp]
\centering
\begin{tabular}{|l|c|c|c|c|}
\hline
\textbf{Class} & \textbf{Precision} & \textbf{Recall} & \textbf{F1 Score} & \textbf{Support} \\
\hline
Real & 0.9838 & 0.9755 & 0.9796 & 52,765 \\
Fake & 0.9889 & 0.9927 & 0.9908 & 158,669 \\
\hline
\end{tabular}
\caption{\label{tab:wangModClassWise}Class-wise performance breakdown for the modified \texttt{Wang2020\_128} spatial expert.}
\end{table}

\subsubsection*{Modified Frequency Expert (Wolter2021\_128) Evaluation}
As shown in Table \ref{tab:waveletsModQuanti} and Table \ref{tab:modWaveleteClasswise}, the standalone \texttt{Wolter2021\_128} wavelet expert achieves an overall accuracy of 95.42\%, a ROC-AUC of 0.9864, an Average Precision of 0.9787, an EER of 5.92\%, and a macro F1-score of 0.9453. While the wavelet expert achieves high sensitivity on synthetic images (correctly flagging 156,702 of 158,669 fakes; Fake recall = 0.9876), it misclassifies 6,295 of 52,765 real images as synthetic (Real recall = 0.8807). This elevated false-positive rate reflects the vulnerability of purely spectral detectors to sharp, natural high-frequency edges in authentic facial photographs.

\begin{table}[htbp]
\centering
\begin{tabular}{|l|c|}
\hline
\textbf{Metric} & \textbf{Value} \\
\hline
Accuracy & 95.42\% \\
ROC AUC & 0.9864 \\
Average Precision & 0.9787 \\
Equal Error Rate (EER) & 5.92\% \\
F1 Score (Macro) & 0.9453 \\
\hline
\end{tabular}
\caption{\label{tab:waveletsModQuanti}Quantitative evaluation of the modified \texttt{Wolter2021\_128} wavelet-packet frequency expert on the held-out validation partition ($n=211{,}434$).}
\end{table}

\begin{table}[htbp]
\centering
\begin{tabular}{|l|c|c|c|c|}
\hline
\textbf{Class} & \textbf{Precision} & \textbf{Recall} & \textbf{F1 Score} & \textbf{Support} \\
\hline
Real & 0.9699 & 0.8807 & 0.9232 & 52{,}765 \\
Fake & 0.9479 & 0.9876 & 0.9673 & 158{,}669 \\
\hline
\end{tabular}
\caption{\label{tab:modWaveleteClasswise}Class-wise performance breakdown for the modified \texttt{Wolter2021\_128} frequency expert.}
\end{table}

\subsubsection*{Bidirectional Multi-Head Attention Fusion (MHA\_128) Evaluation}
\label{subsec:mha_eval}
Table \ref{tab:mhaQuantitative} and Table \ref{tab:mhaCalssWise} present the held-out validation performance of the proposed \texttt{MHA\_128} fusion model. Integrating the spatial and frequency embeddings via bidirectional cross-attention elevates overall accuracy to 98.82\%, increases ROC-AUC to 0.9992, raises Average Precision to 0.9984, and reduces the Equal Error Rate to 1.27\% (with a macro F1-score of 0.9862). Out of 52,765 real validation images, 51,288 are correctly classified (Real precision = 0.9900, recall = 0.9720, F1 = 0.9809), and out of 158,669 synthetic images, 157,955 are correctly detected with only 714 false negatives (Fake precision = 0.9873, recall = 0.9955, F1 = 0.9914). Crucially, dynamic cross-attention mitigates the false-positive bias of the standalone wavelet branch, improving Real-class recall by $+9.13$ percentage points (from 88.07\% to 97.20\%) while simultaneously reducing fake misses by 38.3\% relative to the standalone spatial expert (714 vs.\ 1,158 false negatives).

\begin{table}[htbp]
\centering
\begin{tabular}{|l|c|}
\hline
\textbf{Metric} & \textbf{Value} \\
\hline
Accuracy & 98.82\% \\
ROC AUC & 0.9992 \\
Average Precision & 0.9984 \\
Equal Error Rate (EER) & 1.27\% \\
F1 Score (Macro) & 0.9862 \\
\hline
\end{tabular}
\caption{\label{tab:mhaQuantitative}Quantitative evaluation of the proposed \texttt{MHA\_128} fusion model on the held-out validation partition ($n=211{,}434$).}
\end{table}

\begin{table}[htbp]
\centering
\begin{tabular}{|l|c|c|c|c|}
\hline
\textbf{Class} & \textbf{Precision} & \textbf{Recall} & \textbf{F1 Score} & \textbf{Support} \\
\hline
Real & 0.9900 & 0.9720 & 0.9809 & 52{,}765 \\
Fake & 0.9873 & 0.9955 & 0.9914 & 158{,}669 \\
\hline
\end{tabular}
\caption{\label{tab:mhaCalssWise}Class-wise performance breakdown for the proposed \texttt{MHA\_128} fusion model.}
\end{table}
\newpage
\section*{Feature Space Separability and Explainability Analysis}
The t-SNE projections in Figure \ref{fig:fig7_t_SNE_results} visualize the 128-dimensional latent representations learned by the spatial expert, the wavelet expert, and the fused MHA model across authentic images and five synthetic categories (ADM, DDPM, DDIM, DiffSwap, and GAN manipulations). While the standalone wavelet stream exhibits partial cluster overlap between authentic faces and diffusion samples, the fused MHA representation forms tightly compacted, well-separated clusters across both GAN and diffusion generators, visually corroborating the reduction in EER to 1.27\%.

To inspect spatial attribution, Figure \ref{fig:fig8_GradCam_Fake_Image_DiffusionModel_Detected_ADM} presents Grad-CAM activation maps on synthetic diffusion-generated facial imagery (ADM). The heatmaps demonstrate that the detector concentrates high activation weights around periocular regions, mouth contours, and high-frequency skin-texture boundaries where iterative reverse-diffusion denoising leaves subtle local variance anomalies.

\begin{figure}[htbp]
\centering
\includegraphics[width=\linewidth]{fig7_t_SNE_results.jpg}
\caption{t-SNE visualizations of the 128-dimensional latent embeddings across authentic (Real) and synthetic generator families (ADM, DDPM, DDIM, DiffSwap, and GAN) for the spatial expert, wavelet expert, and proposed MHA fusion model.}
\label{fig:fig7_t_SNE_results}
\end{figure}

\begin{figure}[htbp]
\centering
\includegraphics[width=\linewidth]{fig8_GradCam_Fake_Image_DiffusionModel_Detected_ADM}
\caption{Grad-CAM activation visualizations on synthetic diffusion-generated face samples, illustrating model attention focused on periocular, oral, and skin-texture denoising artifacts. The depicted face is an entirely synthetic, AI-generated sample produced by an Ablated Diffusion Model (ADM) and does not represent a real person.}
\label{fig:fig8_GradCam_Fake_Image_DiffusionModel_Detected_ADM}
\end{figure}
\newpage
\section*{Comparative Baseline Evaluation and Ablation Studies}
\subsection*{Comparison with Retrained Baselines and Static Fusion}
\begin{figure}[htbp]
\centering
\includegraphics[width=\linewidth]{fig9_deepfake_detection_models_Acc_ROCAUC_Recall_EER_AP_results.jpg}
\caption{Comparative evaluation of the proposed MHA fusion model against retrained spatial and frequency backbones, static concatenation (\texttt{Fusion\_128}), and the Xception baseline across Accuracy, F1-Score, ROC-AUC, Average Precision, and Equal Error Rate (EER; lower is better).}
\label{fig:fig12_deepfake_detection_models_Acc_ROCAUC_Recall_EER_AP_results}
\end{figure}
Figure \ref{fig:fig12_deepfake_detection_models_Acc_ROCAUC_Recall_EER_AP_results} and Figure \ref{fig:fig13_deepfake_detection_models_precision_Recall_F1_results} compare the proposed \texttt{MHA\_128} model against the original and modified Wang2020 \cite{ref4} and Wolter2021 \cite{ref8} models, a retrained Xception \cite{ref27, ref28} baseline, and a static concatenation baseline (\texttt{Fusion\_128}, which directly concatenates $[X_s \parallel X_f] \in \mathbb{R}^{256}$ into an MLP without cross-attention). All baselines were trained and evaluated under the identical 80/20 benchmark protocol.

\subsection*{Class-Wise Precision, Recall, and F1 Analysis}
\begin{figure}[htbp]
\centering
\includegraphics[width=\linewidth]{fig10_deepfake_detection_models_precision_Recall_F1_results.jpg}
\caption{Detailed comparison of classification metrics separated by class label. While the baseline Xception model shows strong Recall on the Fake class, the proposed MHA model achieves a superior balance, particularly improving Recall-Real (sensitivity) compared to the individual frequency-domain models.}
\label{fig:fig13_deepfake_detection_models_precision_Recall_F1_results}
\end{figure}
As shown in Figure \ref{fig:fig13_deepfake_detection_models_precision_Recall_F1_results}, standalone frequency models (original Wolter at 94.40\% accuracy and modified Wolter at 95.42\% accuracy) suffer from reduced Recall-Real (88.07\%), generating excessive false alarms on authentic faces. Combining spatial and frequency embeddings resolves this vulnerability, lifting Real-class F1 to 0.9809--0.9810 while outperforming the retrained Xception baseline (98.51\% accuracy, 1.57\% EER).

\subsection*{Ablation Study on Latent Embedding Dimensionality}
To justify the 128-dimensional embedding bottleneck ($D = 128$) used across both domain experts, we conducted an ablation study comparing $D \in \{64, 128, 256\}$ under identical training settings. Compressing the expert projections to $D = 64$ created an overly restrictive information bottleneck that discarded fine-grained cross-band wavelet correlations, reducing validation ROC-AUC to 0.9968 and increasing EER to 1.68\%. Conversely, expanding the representation to $D = 256$ quadrupled the parameter count of the cross-attention projection matrices ($W^Q, W^K, W^V$) and caused early overfitting to training-split compression noise, degrading validation EER to 1.49\%. The $D = 128$ configuration provided the best trade-off between representation capacity, generalization stability, and parameter efficiency.

\subsection*{Quantitative Summary and Statistical Comparison of MHA vs.\ Simple Fusion}
\begin{table}[htbp]
\centering
\caption{Quantitative comparison of the proposed Bidirectional Cross-Attention model (\texttt{MHA\_128}) against static concatenation (\texttt{Fusion\_128}) and retrained baselines on the held-out validation split ($n=211{,}434$). The best value in each column is highlighted in bold.}
\label{tab:model_comparison}
\begin{tabular}{|l|c|c|c|c|}
\hline
\textbf{Model} & \textbf{Accuracy} & \textbf{F1 (Real)} & \textbf{ROC AUC} & \textbf{EER} \\ \hline
Our MHA (\texttt{MHA\_128}, Proposed)     & 98.82\% & 0.9809 & \textbf{0.9992} & \textbf{1.27\%} \\ \hline
Our Simple Fusion (\texttt{Fusion\_128})  & \textbf{98.84\%} & \textbf{0.9810} & 0.9900 & 1.32\% \\ \hline
Our Wang (\texttt{Wang2020\_128})         & 98.73\% & 0.9796 & 0.9987 & 1.54\% \\ \hline
Our Wolter (\texttt{Wolter2021\_128})     & 95.42\% & 0.9232 & 0.9864 & 5.92\% \\ \hline
Wang2020 (Original, Re-trained)           & 98.73\% & 0.9796 & 0.9987 & 1.54\% \\ \hline
Wolter2021 (Original, Re-trained)         & 94.40\% & 0.9062 & 0.9752 & 8.03\% \\ \hline
Xception (Re-trained)                     & 98.51\% & 0.9781 & 0.9984 & 1.57\% \\ \hline
\end{tabular}
\end{table}

\subsubsection*{Why ROC-AUC and EER Outweigh Single-Threshold Accuracy in Forensic Detection}
A careful inspection of Table \ref{tab:model_comparison} shows that static concatenation (\texttt{Fusion\_128}) achieves a marginally higher raw accuracy than \texttt{MHA\_128} ($98.84\%$ vs.\ $98.82\%$, a difference of $0.02\%$ or $\sim 42$ samples out of $211{,}434$), and a $0.0001$ difference in Real F1 ($0.9810$ vs.\ $0.9809$). However, raw accuracy and F1 are evaluated at a single fixed operating threshold ($\tau = 0.5$) on an imbalanced validation set (where synthetic images outnumber authentic images by 3:1). At a fixed $\tau = 0.5$ threshold, minor logit shifts on majority-class boundary samples can slightly raise aggregate accuracy even when tail errors worsen.

In operational multimedia forensics, detectors must be calibrated across varying false-alarm tolerances (e.g., high-precision legal verification vs.\ high-recall platform screening). Consequently, threshold-independent ranking metrics---specifically \textbf{ROC-AUC} and \textbf{Equal Error Rate (EER)}---provide a more rigorous measure of decision reliability. As shown in Table \ref{tab:model_comparison}, static concatenation (\texttt{Fusion\_128}) suffers a notable drop in ROC-AUC to $0.9900$ (even lower than the standalone spatial expert's $0.9987$) because unattenuated noisy wavelet features corrupt extreme logit confidence scores when concatenated statically. In contrast, \texttt{MHA\_128} achieves an ROC-AUC of \textbf{0.9992}---representing a \textbf{92.0\% relative reduction in the area above the ROC curve} ($1 - \text{AUC}$ reduced from $0.0100$ in \texttt{Fusion\_128} to $0.0008$ in \texttt{MHA\_128})---while lowering the EER from $1.32\%$ to \textbf{1.27\%}.

\subsubsection*{Multi-Run Statistical Stability and Hypothesis Testing}
To verify that the threshold-independent advantage of \texttt{MHA\_128} over \texttt{Fusion\_128} is statistically stable and not an artifact of random seed initialization, both fusion models were trained and evaluated across $N = 4$ independent random seeds ($\text{seeds} \in \{42, 100, 2024, 3407\}$). Table \ref{tab:statistical_variance} reports the sample mean, standard deviation (SD), and 95\% confidence intervals ($\text{CI}_{95\%} = \bar{x} \pm t_{0.025, 3} \cdot \frac{s}{\sqrt{4}}$, with $t_{0.025, 3} = 3.182$), alongside two-sided paired $t$-test $p$-values.

\begin{table}[htbp]
\centering
\begin{tabular}{|l|c|c|c|}
\hline
\textbf{Metric ($N=4$ Independent Runs)} & \textbf{Simple Fusion (\texttt{Fusion\_128})} & \textbf{Proposed MHA (\texttt{MHA\_128})} & \textbf{Paired $t$-Test ($p$-value)} \\ \hline
Accuracy (\%, Mean $\pm$ SD) & $\mathbf{98.84 \pm 0.03}$ & $98.82 \pm 0.02$ & $p = 0.284$ (Not Sig.) \\
Accuracy 95\% Confidence Interval & $[98.79\%, 98.89\%]$ & $[98.79\%, 98.85\%]$ & Overlapping CIs \\ \hline
ROC-AUC (Mean $\pm$ SD) & $0.9900 \pm 0.0009$ & $\mathbf{0.9992 \pm 0.0004}$ & $p = 0.0004$ ($p < 0.001$) \\
ROC-AUC 95\% Confidence Interval & $[0.9886, 0.9914]$ & $[0.9986, 0.9998]$ & Non-overlapping CIs \\ \hline
EER (\%, Mean $\pm$ SD) & $1.32 \pm 0.04$ & $\mathbf{1.27 \pm 0.02}$ & $p = 0.0082$ ($p < 0.01$) \\
EER 95\% Confidence Interval & $[1.26\%, 1.38\%]$ & $[1.24\%, 1.30\%]$ & Statistically Significant \\ \hline
Macro F1-Score (Mean $\pm$ SD) & $0.9861 \pm 0.0004$ & $\mathbf{0.9862 \pm 0.0003}$ & $p = 0.612$ (Not Sig.) \\ \hline
\end{tabular}
\caption{\label{tab:statistical_variance}Multi-run statistical comparison ($N=4$ independent seeds) between static concatenation (\texttt{Fusion\_128}) and the proposed bidirectional Multi-Head Attention model (\texttt{MHA\_128}), reporting Mean $\pm$ SD, 95\% Confidence Intervals, and two-sided paired $t$-test significance levels.}
\end{table}

As confirmed in Table \ref{tab:statistical_variance}, the $0.02\%$ difference in single-threshold Accuracy ($p = 0.284$) and the $0.0001$ difference in Macro F1 ($p = 0.612$) between \texttt{Fusion\_128} and \texttt{MHA\_128} are statistically indistinguishable, with substantially overlapping 95\% confidence intervals. Conversely, the improvements achieved by \texttt{MHA\_128} in ROC-AUC ($0.9992 \pm 0.0004$ vs.\ $0.9900 \pm 0.0009$, $p = 0.0004$, non-overlapping 95\% CIs) and EER ($1.27\% \pm 0.02\%$ vs.\ $1.32\% \pm 0.04\%$, $p = 0.0082$) are statistically significant, confirming that bidirectional cross-attention provides superior ranking calibration across operating thresholds.

\subsection*{Modularity Verification with Modern Backbones: ConvNeXt Benchmark}
To verify that the proposed MHA fusion module functions as a backbone-agnostic meta-architecture rather than depending exclusively on ResNet-50, we conducted an additional modularity experiment using Meta's modern ConvNeXt architecture. Replacing \texttt{Wang2020\_128} with a modified \texttt{ConvNeXt\_128} spatial expert yielded a standalone validation accuracy of 98.13\% on the 1.01M-image benchmark. When fused with the frozen \texttt{Wolter2021\_128} frequency expert via the bidirectional MHA module, validation accuracy improved to 98.50\%, confirming that cross-domain attention consistently enhances modern spatial backbones.

\subsection*{Robustness Analysis and Adversarial Stress Testing}
To evaluate model behavior under signal perturbations, we distinguish between \textbf{training-integrated degradation resilience} and \textbf{inference-time adversarial stress testing}. During training, dynamic Gaussian blur ($\sigma \in [0.0, 3.0]$) and lossy JPEG compression ($Q \in \{30, 75\}$) were applied, ensuring that the held-out validation metrics (98.82\% accuracy, 0.9992 AUC) inherently reflect resilience to moderate compression and smoothing within the tested benchmark.

To explicitly test vulnerability to gradient-based evasion at inference time, the frozen pipeline was subjected to a white-box Fast Gradient Sign Method (FGSM) attack ($\epsilon = 8/255$ under $L_\infty$ norm) across the validation partition (Table \ref{tab:fgsm_robustness}). Under white-box FGSM perturbation, \texttt{MHA\_128} retains 90.70\% accuracy and 0.8920 ROC-AUC (with EER rising to 16.26\%), supported by the log-scaled Haar wavelet stream which attenuates uniform spatial pixel perturbations.

\begin{table}[htbp]
\centering
\begin{tabular}{|l|c|c|c|}
\hline
\textbf{Evaluation Setting} & \textbf{Accuracy} & \textbf{ROC AUC} & \textbf{EER} \\ \hline
Unperturbed Validation Baseline ($Q=75$) & 0.9882 & 0.9992 & 0.0127 \\ \hline
White-Box FGSM Adversarial Attack ($\epsilon = 8/255$) & 0.9070 & 0.8920 & 0.1626 \\ \hline
\end{tabular}
\caption{Inference-time robustness evaluation comparing baseline validation performance against white-box FGSM adversarial perturbation.}
\label{tab:fgsm_robustness}
\end{table}

We emphasize that our current robustness evaluation is strictly delimited to training-time blur/JPEG augmentations and single-step FGSM perturbations. Multi-step iterative attacks such as Projected Gradient Descent (PGD), multi-stage social-media recompression cascades, severe downscaling, and photometric exposure/dynamic-range shifts were not empirically benchmarked in this study. In real-world online distribution, facial media frequently undergo extreme illumination, exposure clipping, and Low-Dynamic-Range (LDR) to High-Dynamic-Range (HDR) tone-mapping transformations. Recent image reconstruction studies---such as guided LDR generation for single-image HDR reconstruction using exposure priors from multi-exposure images \cite{ref38}---demonstrate that exposure-induced non-linearities and dynamic range compression substantially alter both spatial luminance gradients and high-frequency wavelet coefficients. Incorporating exposure-prior modeling and iterative adversarial training represents an important direction for hardening multi-domain detectors against real-world photometric and adversarial shifts.

\section*{Discussion, Limitations, and Conclusion}
\subsection*{Discussion of Multi-Domain Interaction and Interpretability}
Our empirical findings highlight an important distinction between single-threshold classification accuracy and threshold-independent forensic reliability. While static feature concatenation (\texttt{Fusion\_128}) achieves comparable accuracy at a fixed 0.5 threshold, it applies fixed linear weights that cannot adapt when one modality produces high-confidence false alarms---resulting in degraded tail ranking (0.9900 ROC-AUC). By introducing four-head bidirectional cross-attention with residual feed-forward refinement, \texttt{MHA\_128} dynamically reweights spatial and spectral subspaces on a per-sample basis, achieving a 0.9992 ROC-AUC and 1.27\% EER while keeping total pipeline latency to 13.1\,ms per image.

\subsection*{Limitations and Challenges in Unseen Cross-Domain Generalization}
To maintain rigorous scientific transparency, we explicitly delineate the scope and limitations of the present study:
\begin{enumerate}
    \item \textbf{Within-Benchmark vs.\ External Cross-Dataset Generalization:} Although our evaluation uses strict video-level and identity-level splits across 1,016,855 images spanning multiple GAN and diffusion generators, both the training (80\%) and validation (20\%) partitions are drawn from the same curated multi-source collection (CelebA, FaceForensics++, FakeAVCeleb, Celeb-DF, and Zenodo DiffFace). Consequently, these results establish strong \textbf{within-benchmark multi-generator discrimination}, which must be clearly distinguished from \textbf{zero-shot external cross-dataset generalization}. In extensive preliminary cross-dataset experiments conducted across external unseen distributions (such as DFDC \cite{ref39}, WildDeepfake \cite{ref40}, and Deepfake-Eval-2024 \cite{ref15}) using diverse spatial, frequency, and fusion backbones, we observed that purely supervised detectors consistently suffer severe accuracy degradation under zero-shot domain shifts. This phenomenon arises because external in-the-wild datasets introduce unaligned camera ISP pipelines, unseen video codecs (e.g., variable-bitrate H.264/H.265 social-media recompression), mismatched facial crop scales, and novel generator architectures whose spectral baselines shift the frozen 128-D embeddings outside the learned attention manifold. Addressing zero-shot cross-dataset generalization therefore requires pairing multi-domain attention fusion with test-time domain adaptation, contrastive self-supervised pretraining, or leave-one-generator-out domain generalization protocols in future work.
    \item \textbf{Scope of Baseline Comparisons:} Our controlled experimental comparisons focused on retrained component backbones (Wang2020 \cite{ref4}, Wolter2021 \cite{ref8}), Xception \cite{ref28}, ConvNeXt, and internal fusion ablations under an identical preprocessing pipeline. Direct retraining and benchmarking against recent specialized generalization frameworks (such as UCF \cite{ref31} or forgery semantics decoupling \cite{ref29}) under leave-one-generator-out protocols remain important future extensions.
    \item \textbf{Adversarial, Photometric, and Fairness Scope:} As noted in our robustness and ethics analyses, evaluation against iterative white-box attacks (e.g., PGD, Carlini-Wagner), multi-generation social-media recompression, dynamic range / exposure variations \cite{ref38}, and demographically stratified subgroup fairness audits on fully annotated datasets remain necessary steps before clinical or legal deployment.
\end{enumerate}

\subsection*{Conclusion}
This study presented a modular multi-domain deepfake detection framework that integrates 128-dimensional spatial RGB embeddings (\texttt{Wang2020\_128}) and three-level Haar wavelet packet embeddings (\texttt{Wolter2021\_128}) via a bidirectional Multi-Head Attention (\texttt{MHA\_128}) fusion architecture. Evaluated on a 1,016,855-image multi-generator benchmark under strict identity- and video-level splits, the proposed MHA model achieves 98.82\% accuracy, 0.9992 ROC-AUC, and 1.27\% EER. Multi-seed statistical testing confirms that bidirectional cross-attention significantly improves threshold-independent ranking calibration ($p < 0.01$) over static feature concatenation while eliminating the high false-positive rate of standalone frequency analysis. Combined with its low computational overhead (1.73M parameters in the wavelet expert and 0.44M in the attention head) and verified modularity with Xception and ConvNeXt backbones, the proposed framework provides an interpretable, extensible foundation for multi-domain synthetic media forensics.

\section*{Declarations}
\subsection*{Funding}
The authors declare that no funds, grants, or other support were received during the preparation of this manuscript. Open access funding was provided by the Science, Technology \& Innovation Funding Authority (STDF) in cooperation with the Egyptian Knowledge Bank (EKB).

\subsection*{Ethics Approval and Consent to Participate}
Not applicable. All datasets utilized in this study (CelebA, FaceForensics++, FakeAVCeleb, Celeb-DF, and Zenodo DiffFace repositories) are publicly available academic benchmarks accessed in accordance with their respective research licenses. No human participants were recruited and no new personally identifiable data were collected.

\subsection*{Consent for Publication}
All sample facial images displayed in the manuscript figures (specifically Figures 1, 2, 3, and 8) are 100\% synthetic, non-existent AI-generated faces produced by generative models (StyleGAN via \texttt{ThisPersonDoesNotExist.com} and Ablated Diffusion Models) and do not depict any real or recognizable human individual; therefore, human subject consent for image publication is not required.

\subsection*{Data Availability}
The datasets analyzed during the current study are publicly available in the Zenodo repository (\url{https://doi.org/10.5281/zenodo.10865300} and \url{https://doi.org/10.5281/zenodo.10865065}) and through the official academic releases of FaceForensics++ \cite{ref12}, FakeAVCeleb \cite{ref11}, CelebA \cite{ref36}, and Celeb-DF \cite{ref37}.

\subsection*{Code Availability}
The complete PyTorch source code, training configurations, and evaluation scripts supporting the findings of this study are publicly archived on Zenodo at \url{https://doi.org/10.5281/zenodo.21323779} \cite{ref26}.

\subsection*{Competing Interests}
The authors declare no competing interests.

\subsection*{Author Contributions}
Abdelrahman Hassan conceived the study, designed the multi-domain MHA architecture, implemented the codebase, conducted the experiments, and wrote the manuscript. Ahmed Hassan and Hossam Abdelmunim supervised the research, validated the methodological design and statistical analyses, and reviewed and edited the manuscript. All authors read and approved the final manuscript.

\begin{thebibliography}{99}

\bibitem{ref1} Ho, J., Jain, A. \& Abbeel, P. Denoising diffusion probabilistic models. \textit{Advances in Neural Information Processing Systems} \textbf{33}, 6840--6851 (2020).

\bibitem{ref2} Song, J., Meng, C. \& Ermon, S. Denoising diffusion implicit models. In \textit{International Conference on Learning Representations (ICLR)} (2021). \url{https://arxiv.org/abs/2010.02502}.

\bibitem{ref3} Dhariwal, P. \& Nichol, A. Diffusion models beat GANs on image synthesis. \textit{Advances in Neural Information Processing Systems} \textbf{34}, 8780--8794 (2021).

\bibitem{ref4} Wang, S.-Y., Wang, O., Zhang, R., Owens, A. \& Efros, A. A. CNN-generated images are surprisingly easy to spot... for now. In \textit{Proceedings of the IEEE/CVF Conference on Computer Vision and Pattern Recognition (CVPR)}, 8695--8704 (2020).

\bibitem{ref5} Ricker, J., Damm, S., Holz, T. \& Fischer, A. Towards the detection of diffusion model deepfakes. In \textit{Proceedings of the 19th International Joint Conference on Computer Vision, Imaging and Computer Graphics Theory and Applications (VISIGRAPP)}, 447--458 (2024). \url{https://arxiv.org/abs/2210.14571}.

\bibitem{ref6} Lin, H., Luo, W., Wei, K. \& Liu, M. Improved Xception with dual attention mechanism and feature fusion for face forgery detection. \textit{IEEE Signal Processing Letters} \textbf{29}, 1639--1643 (2022).

\bibitem{ref7} Yang, S., Guo, H., Hu, S., Zhu, B., Fu, Y., Lyu, S., Wu, X. \& Wang, X. CrossDF: Improving cross-domain deepfake detection with deep information decomposition. \textit{arXiv preprint arXiv:2310.00359} (2023).

\bibitem{ref8} Wolter, M., Blanke, F., Heese, R. \& Garcke, J. Wavelet-packets for deepfake image analysis and detection. \textit{Machine Learning} \textbf{111}, 4295--4327 (2022).

\bibitem{ref9} Chen, Z. DiffFace - part1 [Data set]. \textit{Zenodo} (2024). \url{https://doi.org/10.5281/zenodo.10865300}.

\bibitem{ref10} Chen, Z. DiffFace - part2 [Data set]. \textit{Zenodo} (2024). \url{https://doi.org/10.5281/zenodo.10865065}.

\bibitem{ref11} Khalid, H., Tariq, S., Kim, M. \& Woo, S. S. FakeAVCeleb: A novel audio-video multimodal deepfake dataset. In \textit{Proceedings of the Neural Information Processing Systems Track on Datasets and Benchmarks} \textbf{1} (2021).

\bibitem{ref12} Rössler, A., Cozzolino, D., Verdoliva, L., Riess, C., Thies, J. \& Nießner, M. FaceForensics++: Learning to detect manipulated facial images. In \textit{Proceedings of the IEEE/CVF International Conference on Computer Vision (ICCV)}, 1--11 (2019).

\bibitem{ref13} Wang, Y., Chen, X., Zhu, J., Chu, W., Tai, Y., Wang, C., Li, J., Wu, Y., Huang, F. \& Ji, R. HifiFace: 3D shape and semantic prior guided high fidelity face swapping. In \textit{Proceedings of the Thirtieth International Joint Conference on Artificial Intelligence (IJCAI)}, 1136--1142 (2021).

\bibitem{ref14} Szegedy, C., Ioffe, S., Vanhoucke, V. \& Alemi, A. A. Inception-v4, inception-resnet and the impact of residual connections on learning. In \textit{Proceedings of the AAAI Conference on Artificial Intelligence} \textbf{31}, 4278--4284 (2017).

\bibitem{ref15} Chandra, N. A. et al. Deepfake-Eval-2024: A multi-modal in-the-wild benchmark of deepfakes circulated in 2024. \textit{arXiv preprint arXiv:2503.02857} (2025).

\bibitem{ref16} Zhao, W., Rao, Y., Shi, W., Liu, Z., Zhou, J. \& Lu, J. DiffSwap: High-fidelity and controllable face swapping via 3D-aware masked diffusion. In \textit{Proceedings of the IEEE/CVF Conference on Computer Vision and Pattern Recognition (CVPR)}, 8568--8577 (2023).

\bibitem{ref17} Zhu, Y., Li, Q., Wang, J., Xu, C. \& Sun, Z. One shot face swapping on megapixels. In \textit{Proceedings of the IEEE/CVF Conference on Computer Vision and Pattern Recognition (CVPR)}, 4834--4844 (2021).

\bibitem{ref18} Thies, J., Zollhöfer, M., Stamminger, M., Theobalt, C. \& Nießner, M. Face2Face: Real-time face capture and reenactment of RGB videos. In \textit{Proceedings of the IEEE Conference on Computer Vision and Pattern Recognition (CVPR)}, 2387--2395 (2016).

\bibitem{ref19} Li, M., Yuan, G., Wang, C., Liu, Z., Zhang, Y., Nie, Y., Wang, J. \& Xu, D. E4S: Fine-grained face swapping via editing with regional GAN inversion. \textit{arXiv preprint arXiv:2310.15081} (2023).

\bibitem{ref20} Li, L., Bao, J., Yang, H., Chen, D. \& Wen, F. FaceShifter: Towards high fidelity and occlusion aware face swapping. \textit{arXiv preprint arXiv:1912.13457} (2019).

\bibitem{ref21} Kim, K., Kim, Y., Cho, S., Seo, J., Nam, J., Lee, K., Kim, S. \& Lee, K. DiffFace: Diffusion-based face swapping with facial guidance. \textit{Pattern Recognition} \textbf{163}, 111451 (2025).

\bibitem{ref22} Ye, F., Hua, M., Zhang, P., Li, X., Sun, Q., Zhao, S., He, Q. \& Wu, X. DreamID: High-fidelity and fast diffusion-based face swapping via triplet ID group learning. \textit{arXiv preprint arXiv:2504.14509} (2025).

\bibitem{ref23} Shao, H., Wang, S., Zhou, Y., Song, G., He, D., Qin, S., Zong, Z., Ma, B., Liu, Y. \& Li, H. VividFace: A diffusion-based hybrid framework for high-fidelity video face swapping. \textit{arXiv preprint arXiv:2412.11279} (2024).

\bibitem{ref24} Afchar, D., Nozick, V., Yamagishi, J. \& Echizen, I. MesoNet: A compact facial video forgery detection network. In \textit{2018 IEEE International Workshop on Information Forensics and Security (WIFS)}, 1--7 (2018).

\bibitem{ref25} Wang, S.-Y. CNNDetection GitHub Repository. (2020). \url{https://github.com/PeterWang512/CNNDetection.git}.

\bibitem{ref26} Hassan, A. CNNDetection Retraining and Fusion Framework. \textit{Zenodo} (2026). \url{https://doi.org/10.5281/zenodo.21323779}.

\bibitem{ref27} Standley, T. Xception PyTorch GitHub Repository. (2017). \url{https://github.com/tstandley/Xception-PyTorch.git}.

\bibitem{ref28} Chollet, F. Xception: Deep learning with depthwise separable convolutions. In \textit{Proceedings of the IEEE Conference on Computer Vision and Pattern Recognition (CVPR)}, 1251--1258 (2017).

\bibitem{ref29} Ye, W., He, X. \& Ding, F. Decoupling forgery semantics for generalizable deepfake detection. \textit{arXiv preprint arXiv:2406.09739} (2024).

\bibitem{ref30} Soudy, A. H., Sayed, O., Tag-Elser, H., Ragab, R., Mohsen, S., Mostafa, T. A. H., Abohany, A. A. \& Slim, S. O. Deepfake detection using convolutional vision transformers and convolutional neural networks. \textit{Neural Computing and Applications} \textbf{36}, 19759--19775 (2024).

\bibitem{ref31} Yan, Z., Zhang, Y., Fan, Y. \& Wu, B. UCF: Uncovering common features for generalizable deepfake detection. In \textit{Proceedings of the IEEE/CVF International Conference on Computer Vision (ICCV)}, 22412--22423 (2023).

\bibitem{ref32} Ye, C., Chen, W., Li, J., Zhang, L. \& Mao, Z. Dual-path Collaborative Generation Network for Emotional Video Captioning. In \textit{Proceedings of the 32nd ACM International Conference on Multimedia (ACM MM)}, 496--505 (2024).

\bibitem{ref33} Ye, C., Chen, W., Hu, B., Zhang, L., Zhang, Y. \& Mao, Z. Improving Video Summarization by Exploring the Coherence between Corresponding Captions. \textit{IEEE Transactions on Image Processing} \textbf{34}, 5369--5384 (2025).

\bibitem{ref34} Ye, C., Chen, W., Song, P., Liu, X., Zhang, L. \& Mao, Z. Multi-round Mutual Emotion-Cause Pair Extraction for Emotion-Attributed Video Captioning. In \textit{Proceedings of the 33rd ACM International Conference on Multimedia (ACM MM)}, 3320--3329 (2025).

\bibitem{ref35} Chen, W., Ye, C., Song, P., Zhang, L., Zhang, Y. \& Mao, Z. Subjective-objective Emotion Correlated Generation Network for Subjective Video Captioning. \textit{IEEE Transactions on Image Processing} \textbf{35}, 540--555 (2026).

\bibitem{ref36} Liu, Z., Luo, P., Wang, X. \& Tang, X. Deep learning face attributes in the wild. In \textit{Proceedings of the IEEE International Conference on Computer Vision (ICCV)}, 3730--3738 (2015).

\bibitem{ref37} Li, Y., Yang, X., Sun, P., Qi, H. \& Lyu, S. Celeb-DF: A large-scale challenging dataset for deepfake forensics. In \textit{Proceedings of the IEEE/CVF Conference on Computer Vision and Pattern Recognition (CVPR)}, 3207--3216 (2020).

\bibitem{ref38} Zheng, B., Guo, S., Zhang, Q., Wang, H. et al. Guided LDR Generation for Single-Image HDR Reconstruction Using Exposure Priors From Multi-Exposure Images. \textit{IEEE Transactions on Circuits and Systems for Video Technology} \textbf{36}, 15478--15491 (2026). \url{https://doi.org/10.1109/TCSVT.2026.3709976}.

\bibitem{ref39} Dolhansky, B. et al. The DeepFake Detection Challenge (DFDC) dataset. \textit{arXiv preprint arXiv:2006.07397} (2020).

\bibitem{ref40} Zi, B., Chang, M., Chen, J., Ma, X. \& Jiang, Y.-G. WildDeepfake: A challenging real-world dataset for deepfake detection. In \textit{Proceedings of the 28th ACM International Conference on Multimedia (ACM MM)}, 2382--2390 (2020).

\end{thebibliography}

\section*{Reproducibility and Environment Specifications}

To ensure the reproducibility of our experimental results and facilitate independent verification, we provide complete software and hardware environment specifications, alongside a structured reproduction protocol.

\subsection*{Hardware Environment}
All stages of the experimental pipeline were conducted across three clearly separated hardware configurations, validating both resource-constrained prototyping and homogeneous synchronized multi-GPU training:
\begin{enumerate}
    \item \textbf{Prototyping \& Development Stage:} A local workstation equipped with a single NVIDIA GeForce RTX 3060 GPU (12\,GB GDDR6 VRAM), an Intel Core i5-12400 CPU (6 cores, 12 threads), 32\,GB DDR4 RAM, and 2\,TB HDD storage, used exclusively for initial pipeline debugging and hyper-parameter verification.
    \item \textbf{Production Training \& Main Evaluation Stage:} A homogeneous cloud-based multi-GPU node consisting of two identical NVIDIA Tensor Core T4 GPUs (16\,GB GDDR6 VRAM each) interconnected via PCIe Gen3 (the RTX 3060 was not combined with T4 GPUs during distributed training). Gradient synchronization and data distribution across the two T4 GPUs were managed using PyTorch's native \texttt{DistributedDataParallel} (DDP) API with the \texttt{nccl} backend and fixed random seeds, ensuring deterministic gradient averaging across parallel mini-batches.
    \item \textbf{Cross-Hardware Verification Stage:} An isolated cloud instance featuring a single NVIDIA Ada Lovelace L4 GPU (24\,GB GDDR6 VRAM). This was used to verify runtime metrics and confirm the hardware-agnostic stability of the trained checkpoints, which exhibited negligible metric variance ($< 0.05\%$) compared to the homogeneous dual-T4 setup.
\end{enumerate}

\subsection*{Software Environment and Dependencies}
The software stack was standardized to prevent dependency-induced discrepancies. The complete environment is defined by the following specifications:
\begin{itemize}
    \item \textbf{Operating System:} Windows 10/11 and Linux (Ubuntu 20.04 LTS or newer).
    \item \textbf{Python Version:} $\geq$ 3.11 (specifically tested and executed on Python 3.11.0).
    \item \textbf{CUDA Support:} CUDA Runtime version 12.1.
    \item \textbf{Core Libraries:} PyTorch v2.5.1 and torchvision v0.20.1, both compiled with CUDA 12.1 support.
\end{itemize}
The exact versions of auxiliary Python libraries used for data loading, preprocessing, and metric evaluation are detailed in Table~\ref{tab:software_env}.

\begin{table}[ht]
\centering
\caption{Complete software dependency and package version specifications.}
\label{tab:software_env}
\begin{tabular}{|l|c|l|}
\hline
\textbf{Library Name} & \textbf{Minimum Version} & \textbf{Primary Purpose in Pipeline} \\ \hline
\texttt{torch} & 2.5.1 & Core deep learning framework and tensor operations \\ \hline
\texttt{torchvision} & 0.20.1 & Data loading utilities and image transformations \\ \hline
\texttt{numpy} & 2.4.4 & Multi-dimensional numerical arrays and calculations \\ \hline
\texttt{pytorch-wavelets} & 1.3.0 & GPU-accelerated Discrete Wavelet Transforms \\ \hline
\texttt{pywavelets} & 1.9.0 & CPU-based reference wavelet packet decompositions \\ \hline
\texttt{scikit-learn} & 1.9.0 & Evaluation metric calculations (EER, AUC, F1) \\ \hline
\texttt{opencv-python} & 5.0.0.93 & Real-time image decoding and augmentations \\ \hline
\texttt{pillow} & 12.2.0 & Image processing and basic loading checks \\ \hline
\texttt{matplotlib} & 3.11.0 & Metric visualization (ROC, PR, Confusion Matrix) \\ \hline
\texttt{scipy} & 1.17.1 & Mathematical utilities and signal processing algorithms \\ \hline
\texttt{pyyaml} & 6.0.3 & Parsing configuration files and CLI overrides \\ \hline
\texttt{tensorboard} & 2.21.0 & Logging training progress and metric visualization \\ \hline
\texttt{tqdm} & 4.68.3 & Console-based execution progress tracking \\ \hline
\end{tabular}
\end{table}

\subsection*{Dataset and Preprocessing Specifications}
The training and validation datasets~\cite{ref9,ref10,ref11,ref12,ref36,ref37} follow a standardized two-level hierarchical folder structure compatible with PyTorch's \texttt{ImageFolder} representation:
\begin{itemize}
    \item \textbf{Directory Layout:} Root folders named \texttt{train} and \texttt{val}, each containing two subfolders named \texttt{real} and \texttt{fake}.
    \item \textbf{Image Preprocessing:} Images are dynamically resized to a fixed resolution of $224 \times 224$ pixels. Normalization is performed on input tensors using the ImageNet channel-wise mean ($[0.485, 0.456, 0.406]$) and standard deviation ($[0.229, 0.224, 0.225]$).
\end{itemize}

\subsection*{Step-by-Step Reproduction Protocol}
The source code, training configurations, and evaluation scripts are publicly archived at~\cite{ref26}. To reproduce the experiments:
\begin{enumerate}
    \item \textbf{Cloning and Virtual Environment Setup:}
    Clone the codebase and initialize a Python environment using either the package manager \texttt{uv} (recommended for deterministic package resolution via \texttt{uv.lock}) or standard \texttt{pip}:
    \begin{verbatim}
    git clone <repository_url> && cd "Novel Deepfake Detection"
    
    # Option A: Fast installation via uv
    uv venv && .venv\Scripts\activate
    uv sync
    
    # Option B: Standard installation via pip
    python -m venv .venv && .venv\Scripts\activate
    pip install -r requirements.txt
    \end{verbatim}

    \item \textbf{Phase 1: Base Model Training (128-D Embeddings):}
    Train the backbone models (Wang2020, Wolter Wavelet, Xception, ConvNeXt) on the 80\% training split to output 128-D embeddings:
    \begin{verbatim}
    python train.py --arch Wang2020_128 --dataroot "path/to/train" --val_root "path/to/val"
    python train.py --arch WolterWavelet2021_128 --dataroot "path/to/train" ...
    python train.py --arch Xception_128 --dataroot "path/to/train" ...
    python train.py --arch ConvNeXt_128 --dataroot "path/to/train" ...
    \end{verbatim}

    \item \textbf{Phase 2: Fusion and Multi-Head Attention Model Training:}
    Train the fusion networks (e.g., \texttt{MHA\_128} and \texttt{Fusion\_128}) on the 80\% training split by providing paths to the frozen Phase 1 checkpoints:
    \begin{verbatim}
    python train.py --arch MHA_WWXC --dataroot "path/to/train" --val_root "path/to/val" \
      --rgb_model_path "experiments/wang2020_128_run/checkpoints/best.pth" \
      --wavelet_model_path "experiments/wolter_128_run/checkpoints/best.pth" \
      --xception_model_path "experiments/xception_128_run/checkpoints/best.pth" \
      --convnext_model_path "experiments/convnext_128_run/checkpoints/best.pth" \
      --freeze_base_models
    \end{verbatim}

    \item \textbf{Phase 3: Held-Out Model Evaluation and FGSM Stress Testing:}
    Run comparative evaluations on the held-out 20\% validation split and test inference-time robustness under white-box FGSM perturbation:
    \begin{verbatim}
    # Comprehensive evaluation & multi-seed statistical report generation
    python evaluate.py --val_root "path/to/val" --mha_wwxc "path/to/best_mha.pth" ...
    
    # White-box FGSM adversarial stress testing
    python evaluate_robustness.py --val_root "path/to/val" --attack fgsm --mha_wwxc "path/to/best_mha.pth" ...
    \end{verbatim}
\end{enumerate}

\end{document}